\chapter{Geodesic spheres and spherical CMC rigidity}\label{sec:cmc-spaceforms}
The scalar constraint makes geodesic spheres stable even when their unrestricted second variation has a negative direction. We develop this explicitly in space forms and prove the spherical rigidity argument of \cite{BdCE88}, Section 3, pp.~129--131. All hypersurfaces here have dimension $n\geq2$. We use unaveraged $H=\operatorname{tr}(D\nu)$ and the stability convention in Definition \ref{def:cmc-stability-types}.

\begin{proposition}[Geometry of a space-form geodesic sphere]\label{prop:spaceform-sphere-geometry}
In the simply connected complete space form of sectional curvature $c$, let $\Sigma_r$ be a geodesic sphere of radius $r>0$, with $r<\pi/\sqrt c$ when $c>0$. Set
\[
 s_c(r)=\begin{cases}\sin(\sqrt c\,r)/\sqrt c,&c>0,\\ r,&c=0,\\ \sinh(\sqrt{-c}\,r)/\sqrt{-c},&c<0.\end{cases}
\]
With the normal pointing away from the center, its induced metric is $s_c(r)^2g_{S^n}$ and
\[
 A=\frac{s_c'}{s_c}I,\qquad H=n\frac{s_c'}{s_c},\qquad
 q=|A|^2+\operatorname{Ric}(\nu,\nu)=\frac n{s_c^2}.
\]
These formulas make explicit the quantities in \cite{BdCE88}, Lemma 3.4.
\end{proposition}
\begin{proof}
A perpendicular Jacobi field along a unit radial geodesic with initial data $Y(0)=0$, $Y'(0)=v$ satisfies $Y''+cY=0$ after parallel transport. Thus $Y(r)=s_c(r)v$, giving the polar metric $dr^2+s_c(r)^2g_{S^n}$ before the antipode in the positive-curvature case. Taking one half of the radial derivative of the tangential metric gives $A=(s_c'/s_c)I$. Also $\operatorname{Ric}(\nu,\nu)=nc$. The identity $(s_c')^2+cs_c^2=1$ follows either from the displayed formulas or by differentiating it and using $s_c''+cs_c=0$ and the initial data. Therefore $q=n((s_c'/s_c)^2+c)=n/s_c^2$.
\end{proof}

\begin{theorem}[Stability and unrestricted index of geodesic spheres]\label{thm:spaceform-sphere-stability}
For $n\geq2$, every geodesic sphere in Proposition \ref{prop:spaceform-sphere-geometry} is stable for the enclosed-volume constraint and has unrestricted index exactly one. This is the stability assertion of \cite{BdCE88}, Lemma 3.4, together with its index consequence.
\end{theorem}
\begin{proof}
The induced sphere has radius $s_c(r)$, so Proposition \ref{prop:round-sphere-poincare} gives first positive eigenvalue $n/s_c(r)^2=q$. Proposition \ref{prop:cmc-constant-potential-test} proves constrained stability. The constant function gives $Q(1,1)=-q\operatorname{Area}(\Sigma_r)<0$. Proposition \ref{prop:cmc-constrained-index} then bounds the unrestricted index above by one, proving equality. The exact-volume realization proposition identifies the mean-zero calculation with the geometric variation problem.
\end{proof}

\begin{proposition}[Nullity of a geodesic sphere]\label{prop:spaceform-sphere-nullity}
The constrained Jacobi nullspace of $\Sigma_r$ has dimension $n+1$ and consists of the first spherical harmonics for its round induced metric. The unrestricted Jacobi nullspace is the same space.
\end{proposition}
\begin{proof}
The potential is the positive constant $q=n/s_c(r)^2$. By Proposition \ref{prop:cmc-constant-potential-index}, the nullspace is the Laplace eigenspace with eigenvalue $q$. Proposition \ref{prop:spherical-harmonic-h1-expansion}, scaled to the round induced radius $s_c(r)$ as in Corollary \ref{cor:sphere-higher-mode-gap}, identifies every weak eigenfunction with a harmonic degree and gives eigenvalues $\ell(\ell+n-1)/s_c(r)^2$. The value $q$ occurs exactly at $\ell=1$. Example \ref{ex:sphere-low-harmonic-modes} identifies this space as the restrictions of linear functions on the intrinsic sphere model; it has dimension $n+1$ and zero mean. Constants have Laplace eigenvalue zero, so they add no unrestricted Jacobi fields. This proves both nullity statements using the local spectral decomposition. It concerns infinitesimal nullity and does not assert uniqueness of a nonlinear CMC deformation.
\end{proof}

\begin{lemma}[Position and normal equations in the unit sphere]\label{lem:spherical-cmc-height-equations}
Let $\Sigma^n\subset S^{n+1}\subset\mathbb R^{n+2}$ be a smooth two-sided CMC hypersurface with position $x$ and unit normal $\nu$ tangent to $S^{n+1}$. Then
\[
 \Delta x=-nx-H\nu,\qquad \Delta\nu=-H x-|A|^2\nu.
\]
For $g_v=\langle x,v\rangle$, $f_v=\langle\nu,v\rangle$ and constant $v\in\mathbb R^{n+2}$, these imply
\[
 Jg_v=|A|^2g_v-Hf_v,\qquad Jf_v=nf_v-Hg_v,
 \qquad J=\Delta+|A|^2+n.
\]
Compare \cite{BdCE88}, Lemma 3.3, after converting its mean-curvature convention.
\end{lemma}
\begin{proof}
Choose an orthonormal tangent frame geodesic at the point in question. In the ambient Euclidean space,
\[
 D_{e_i}e_j=\nabla^\Sigma_{e_i}e_j-h_{ij}\nu-\delta_{ij}x,
 \qquad D_{e_i}\nu=\sum_jh_{ij}e_j.
\]
Tracing the first formula gives $\Delta x$. Differentiating the second and tracing gives
$\Delta\nu=\nabla H-|A|^2\nu-Hx$, using Codazzi in constant curvature. Since $H$ is constant, the gradient vanishes. Pairing the vector equations with $v$ and adding the potential $|A|^2+n$ gives the claimed Jacobi formulas. The ambient position and hypersurface normal are perpendicular, both of unit length.
\end{proof}

\begin{proposition}[Spherical coordinate tests and the summed defect]\label{prop:spherical-cmc-defect-test}
For a closed hypersurface in Lemma \ref{lem:spherical-cmc-height-equations}, put $D=|A|^2-H^2/n$ and
\[
 u_v=g_v+\frac Hn f_v.
\]
Then $\int_\Sigma u_v=0$ and $Ju_v=Dg_v$. If $v_1,\ldots,v_{n+2}$ is an orthonormal basis of $\mathbb R^{n+2}$, then
\[
 \sum_{i=1}^{n+2}Q(u_{v_i},u_{v_i})=-\int_\Sigma D\,dA.
\]
This is the spherical test family in \cite{BdCE88}, proof of Theorem 1.2, pp.~130--131.
\end{proposition}
\begin{proof}
The position equation gives $u_v=-\Delta g_v/n$, hence zero integral. The two Jacobi equations give
\[
 Ju_v=(|A|^2-H^2/n)g_v+(-H+H)f_v=Dg_v.
\]
Therefore $Q(u_v,u_v)=-\int D(g_v^2+(H/n)f_vg_v)$. Summing over the basis uses $\sum_i g_{v_i}^2=|x|^2=1$ and $\sum_i f_{v_i}g_{v_i}=\langle\nu,x\rangle=0$. These pointwise identities give the integral formula. Each summand is an admissible volume-constrained test, whether or not it vanishes identically.
\end{proof}

\begin{theorem}[Volume-stable CMC boundaries in a round sphere]\label{thm:spherical-cmc-stability-rigidity}
For $n\geq2$, a smooth closed connected embedded two-sided CMC boundary in the unit sphere $S^{n+1}$ is volume-constrained stable if and only if it is a geodesic sphere, including an equator. This is the positive-curvature case of \cite{BdCE88}, Theorem 1.2.
\end{theorem}
\begin{proof}
Set $D=|A|^2-H^2/n$. Under stability, every summand in Proposition \ref{prop:spherical-cmc-defect-test} is nonnegative, whereas their sum is $-\int D\leq0$. It follows that $D=0$ everywhere, hence $A=\kappa I$ with the constant $\kappa=H/n$. In $\mathbb R^{n+2}$,
\[
 D_v(\nu-\kappa x)=Av-\kappa v=0
\]
for every tangent vector $v$. Connectedness therefore gives a constant vector $a=\nu-\kappa x$, with $|a|^2=1+\kappa^2$ and $\langle x,a\rangle=-\kappa$. The intersection of the unit sphere with this affine hyperplane is a connected geodesic sphere: its offset has absolute value $|\kappa|/\sqrt{1+\kappa^2}<1$. The hypersurface image is open in that intersection by equality of tangent dimensions and closed by compactness, hence equals it. The converse is Theorem \ref{thm:spaceform-sphere-stability}. When $\kappa=0$, the hyperplane passes through the origin and the geodesic sphere is an equator. Uniform rescaling gives the corresponding result for every positive ambient constant curvature.
\end{proof}

\begin{example}[An equator separates two notions of minimal stability]\label{ex:equator-constrained-stability}
An equator $S^n\subset S^{n+1}$ is minimal, has $A=0$, and is stable for the enclosed-volume constraint but unstable for unrestricted area variations. Its unrestricted index is one. This example is emphasized in \cite{BdCE88}, Remark 3.5.
\end{example}
\begin{proof}
At geodesic radius $\pi/2$ the space-form formula gives $H=0$ and $q=n$ because the ambient normal Ricci curvature is $n$. Thus $Q(f,f)=\int(|\nabla f|^2-nf^2)$. The sharp spherical Poincare inequality makes this nonnegative for zero-mean $f$, while $Q(1,1)<0$. The index statement follows as above. In particular the minimal equation $H=0$ does not remove the distinction between testing all normal speeds and testing only those that preserve enclosed volume.
\end{proof}

\begin{example}[Product tori have many volume-preserving unstable directions]\label{ex:cmc-product-torus-instability}
Let $p,q\geq1$, $n=p+q$, $r,s>0$, and $r^2+s^2=1$. The product
\[
 \Sigma=S^p(r)\times S^q(s)\subset S^{n+1}
\]
is a closed embedded CMC hypersurface. Its constrained index is at least $n+2$, and its unrestricted index is at least $n+3$. This supplies explicit unstable tests consistent with Theorem \ref{thm:spherical-cmc-stability-rigidity}.
\end{example}
\begin{proof}
At $(x,y)$ take $\nu=((s/r)x,-(r/s)y)$. It is unit and perpendicular both to the position and to the product tangent space. Differentiating gives principal curvatures $s/r$ in the $p$ first-factor directions and $-r/s$ in the $q$ other directions. Hence $H=ps/r-qr/s$ is constant and
\[
 |A|^2+n=\frac p{r^2}+\frac q{s^2}.
\]
Each of the $p+1$ first-factor coordinate functions has zero mean and Laplace eigenvalue $p/r^2$, so its quadratic form is $-(q/s^2)\int f^2$. Each of the $q+1$ second-factor coordinates has eigenvalue $q/s^2$ and quadratic form $-(p/r^2)\int f^2$. Coordinate functions in this combined space are mutually orthogonal, and the second-variation form is diagonal on it, including when the two eigenvalues coincide. Thus they span a negative-definite mean-zero space of dimension $p+q+2=n+2$. The constant function is orthogonal to it and gives another negative direction. The coordinate eigenvalue formula follows directly by tracing the Euclidean second fundamental form of each round factor; no full product spectral decomposition is needed.
\end{proof}
